\section{Setup}\label{sec:08-group-theorems:01-setup:1}

Chapter 7 closed by pointing out the entire payoff of bundling data with proofs in the first place: a theorem proved once about a generic \texttt{Grp : Group G} is inherited for free by \texttt{intGroup}, by \texttt{perm3Group}, and by every group built afterward, with nothing about the specific carrier re-proved at each use site. This chapter cashes in that promise. Rather than fix \texttt{G} to \texttt{Int} or \texttt{Perm3} before proving anything, the setup below fixes an \emph{arbitrary} group once, and every theorem that follows is proved against that one unnamed group, so that whatever gets proved here automatically applies everywhere a \texttt{Group} instance is later supplied.

\begin{lstlisting}
variable {G : Type} (Grp : Group G)
\end{lstlisting}

\texttt{variable} adds \texttt{\{G : Type\} (Grp : Group G)} silently to every following declaration that mentions \texttt{G} or \texttt{Grp}.

\begin{mathreading}

This is the phrase ``Let \(G\) be a group'' opening a section: we fix an arbitrary object \((G, \cdot, e, (-)^{-1})\) of \(\mathbf{Grp}\) and reason generically about it. Everything proved under this \texttt{variable} is a statement quantified over \emph{all} groups. A theorem about the group \(\mathrm{Grp}\) is really \(\forall (G : \mathrm{Type})\,(\mathrm{Grp} : \mathrm{Group}\,G),\ (\ldots)\), so it applies to \(\mathbb{Z}\), to permutation groups, and to every group built later.

\end{mathreading}

\subsection{Sources, quoted}\label{sec:08-group-theorems:01-setup:2}

Formal definitions and citations for this section, gathered here for reference (full entries in the \hyperref[sec:root:bibliography:1]{Bibliography}):

\begin{itemize}
\tightlist
\item
  \textbf{Group theorems proved here.} ``The identity of \(G\) is unique \ldots{} for each \(a \in G\), \(a^{-1}\) is uniquely determined \ldots{} \((a{*}b)^{-1} = (b^{-1}){*}(a^{-1})\)'' (\cite{DummitFoote2003}, §1.1 ``Basic Axioms and Examples,'' pp.~17--18, Proposition 1).
\end{itemize}
